\documentclass[a4paper]{article}

% Basic packages
\usepackage[fontset=fandol]{ctex}
\usepackage{amsmath, amssymb}
\usepackage{graphicx}
\usepackage[margin=2.5cm]{geometry}
\usepackage[most]{tcolorbox}
\usepackage{etoolbox}
\usepackage{listings}
\usepackage{hyperref}
\usepackage{booktabs}
\usepackage{subcaption}
\usepackage{float}
\usepackage{tikz}
\usetikzlibrary{positioning}
\IfFileExists{pgfplots.sty}{
    \usepackage{pgfplots}
    \pgfplotsset{compat=1.18}
}{}

% Highlight boxes
\newtcolorbox{knowledgebox}[1]{
    enhanced, colback=blue!5!white, colframe=blue!75!black, colbacktitle=blue!75!black,
    coltitle=white, fonttitle=\bfseries, title=#1, attach boxed title to top left={yshift=-2mm, xshift=2mm},
    boxrule=1pt, sharp corners
}

\newtcolorbox{importantbox}[1]{
    enhanced, colback=yellow!10!white, colframe=yellow!80!black, colbacktitle=yellow!80!black,
    coltitle=black, fonttitle=\bfseries, title=#1, sharp corners
}

\newtcolorbox{warningbox}[1]{
    enhanced, colback=red!5!white, colframe=red!75!black, colbacktitle=red!75!black,
    coltitle=white, fonttitle=\bfseries, title=#1, sharp corners
}

% Listings
\lstset{
    language=Python,
    basicstyle=\ttfamily\small,
    keywordstyle=\color{blue},
    stringstyle=\color{red!60!black},
    commentstyle=\color{green!60!black},
    breaklines=true,
    frame=single,
    numbers=left,
    numberstyle=\tiny\color{gray},
    captionpos=b,
    extendedchars=false
}

% Front-page metadata
\newcommand{\notetitle}{CS224N Lecture 9: Pretraining}
\newcommand{\noteauthors}{整理：AI Course Notes Contributors；授课：Chris Manning}
\newcommand{\notedate}{2024年}
\newcommand{\videochannel}{Stanford Online}
\newcommand{\videopublishdate}{2024}
\newcommand{\videoduration}{1:18:00}
\newcommand{\videourl}{https://www.youtube.com/watch?v=DGfCRXuNA2w}
\newcommand{\videocoverpath}{}
\newcommand{\repourl}{https://github.com/hqhq1025/ai-course-notes}


% Footer with repo info
\usepackage{fancyhdr}
\pagestyle{fancy}
\fancyhf{}
\fancyfoot[C]{\thepage}
\fancyfoot[R]{\footnotesize\color{gray}\href{\repourl}{AI Course Notes}}
\renewcommand{\headrulewidth}{0pt}
\renewcommand{\footrulewidth}{0pt}

\begin{document}
\setlength{\emergencystretch}{2em}

\begin{titlepage}
\centering
{\Large 课程笔记\par}
\vspace{1.2cm}
{\huge\bfseries \notetitle\par}
\vspace{0.8cm}
{\large \noteauthors\par}
\vspace{0.3cm}
{\large \notedate\par}
\vspace{1.2cm}

\vspace{1.2cm}
\vfill
\begin{tcolorbox}[width=0.9\textwidth, colback=blue!3!white, colframe=blue!55!black, sharp corners]
\textbf{课程版本：} 2024 年中文讲义整理版\par
\textbf{笔记来源：} AI Course Notes Contributors\par
\textbf{本版处理：} 移除课件截图和对应图注，不包含 Stanford 原始课件、图片或视频。\par
\textbf{授权：} 笔记依 CC BY-NC-SA 4.0 分享；完整许可与改动记录见下一页。
\end{tcolorbox}
\end{titlepage}

\begin{center}
\fbox{\begin{minipage}{0.91\textwidth}
\small
\textbf{来源与许可}\\
本中文笔记由 AI Course Notes Contributors 整理，课程版本为 2024；源稿：\url{https://github.com/hqhq1025/ai-course-notes/tree/717e2df65c3d2eee593c171156a30d4f8f6812b9/cs224n/lecture09}。\\
笔记内容依 CC BY-NC-SA 4.0 分享：\url{https://creativecommons.org/licenses/by-nc-sa/4.0/}。\\
本副本的改动：移除 Stanford 原始课件截图与依赖这些截图的图注，移除课程视频封面/链接，并清理 TeX 源稿中的行尾空白后重新编译为可独立阅读的 PDF；同时加入本来源与许可说明。完整许可文本随本文件夹提供。Stanford 课件、课件图片和视频不包含在该授权内；请从对应年份 Stanford 课程页访问原始课件。
\end{minipage}}
\end{center}
\clearpage

\tableofcontents
\newpage

%% ========== Part 1: 子词分词 ==========

\section{子词分词：从有限词表到开放词表}

在课程的前几讲中，我们假设模型拥有一个\textbf{固定的词表} $V$。每个单词被映射到一个唯一的向量表示。然而，这种设计面临一个根本性问题：语言是\textbf{开放的}，新词、拼写变体、形态变化不断出现，固定词表无法覆盖所有可能的输入。


% Raster figure omitted from this study copy.



\subsection{固定词表的问题}

当模型遇到训练时未见过的字符串时，只能将其映射为一个通用的 \texttt{<UNK>}（unknown）标记。这意味着：

\begin{itemize}
    \item 所有未知词共享同一个向量，\textbf{完全丢失了语义信息}
    \item 在形态丰富的语言（如斯瓦希里语、芬兰语、土耳其语）中，一个动词可能有数百种变形，词表膨胀极为严重
    \item 即使是英语，拼写变体、新造词、专有名词也层出不穷
\end{itemize}


% Raster figure omitted from this study copy.

\footnotetext{Slides 第6页。来源：Wiktionary。}

\subsection{Byte Pair Encoding (BPE) 算法}

现代 NLP 的解决方案是\textbf{子词分词}（subword tokenization）。其核心思想是：不以``词''为最小单位，而是以\textbf{频繁共现的字符子串}为基本单元。

\begin{importantbox}{BPE 算法流程}
\begin{enumerate}
    \item 以所有单个字符（加上词尾标记）初始化词表
    \item 统计训练语料中所有相邻字符对的共现频率
    \item 将频率最高的字符对合并为一个新符号，加入词表
    \item 用新符号替换语料中的对应字符对
    \item 重复步骤 2--4，直到词表达到预设大小（通常 30k--50k）
\end{enumerate}
\end{importantbox}


% Raster figure omitted from this study copy.



分词时，算法尽量用\textbf{最少的子词}来表示输入。例如：
\begin{itemize}
    \item 常见词 ``hat'' 直接作为一个 token
    \item 不常见的 ``taaaasty'' 被拆分为 ``ta'' + ``\#\#aa'' + ``\#\#sty'' 三个子词 token
    \item ``transformerify'' 可能被拆分为 ``transformer'' + ``\#\#ify''
\end{itemize}

其中 ``\#\#'' 前缀表示这个子词\textbf{不以空格开头}，即它是前一个子词的延续。


% Raster figure omitted from this study copy.



\begin{knowledgebox}{子词分词的关键性质}
\begin{itemize}
    \item 系统\textbf{不区分}一个 token 是完整的词还是子词片段——它们在模型中的地位完全相同，都是嵌入矩阵中的一个索引
    \item 如果需要获取一个被拆分的词的表示，可以\textbf{取平均}或取最后一个子词的上下文表示
    \item 现代 LLM（GPT、BERT、T5 等）几乎全部使用 BPE 或类似的子词分词方案
    \item 序列长度是 Transformer 的瓶颈，因此分词时倾向于用\textbf{更少更大}的子词
\end{itemize}
\end{knowledgebox}

\begin{warningbox}{子词分词不是``聪明''的语素切分}
BPE 是纯粹基于统计频率的算法，\textbf{不理解语言学意义上的词素}（morpheme）。它可能把 ``unhappiness'' 切分为 ``un'' + ``happiness''（恰好对应语素），也可能切分为 ``unh'' + ``appiness''（无语言学意义）。切分结果完全取决于训练语料的统计分布。
\end{warningbox}

\subsection{本章小结}

子词分词解决了固定词表的 UNK 问题，使模型能够处理任意输入字符串。BPE 通过迭代合并频繁共现字符对来构建词表，在词表大小和序列长度之间取得平衡。这一技术已成为所有现代语言模型的标配。

%% ========== Part 2: 从词嵌入到预训练 ==========

\section{从词嵌入到全模型预训练}

\subsection{分布式假设的两个层次}

课程开篇曾引用 J.R. Firth 的名言：``You shall know a word by the company it keeps.'' 这是\textbf{分布式假设}的经典表述，也是 Word2Vec 的理论基础。

但同一位语言学家还有另一句话：``The complete meaning of a word is always contextual, and no study of meaning apart from a complete context can be taken seriously.'' 这揭示了一个更深层的洞察：词义不仅由其典型的共现伙伴决定，更由其\textbf{具体出现的上下文}决定。


% Raster figure omitted from this study copy.



\begin{importantbox}{Word2Vec 的根本局限}
Word2Vec 为每个词学习\textbf{一个固定的向量}，无论它出现在什么上下文中。对于多义词（如英语的 ``record''：动词``记录'' / 名词``唱片''），这个向量只能是所有义项的\textbf{混合体}，无法在特定语境中区分不同含义。
\end{importantbox}

\subsection{旧范式 vs 新范式}

在预训练出现之前，NLP 的典型工作流程是：

\begin{enumerate}
    \item 用 Word2Vec / GloVe 训练词嵌入（预训练的\textbf{仅仅是嵌入层}）
    \item 在嵌入层之上堆叠一个随机初始化的 LSTM 或 Transformer
    \item 用下游任务（如情感分析、机器翻译）的标注数据从头训练上层网络
\end{enumerate}

新范式的革命性变化在于：\textbf{预训练整个网络}（嵌入层 + Transformer），而不仅仅是词嵌入。


% Raster figure omitted from this study copy.



这带来了三大优势：

\begin{enumerate}
    \item \textbf{上下文表示}：``record'' 在不同语境中获得不同的向量
    \item \textbf{强参数初始化}：所有参数都从一个``好的起点''开始，而非随机初始化
    \item \textbf{语言概率分布}：模型学到了自然语言的生成概率，可以直接用于文本生成
\end{enumerate}

\subsection{预训练的核心思想：重构输入}


% Raster figure omitted from this study copy.



预训练的核心假设是：如果我们遮住文本的一部分，让神经网络来\textbf{重构被遮住的内容}，那么网络为了做好这个任务，必须学会大量关于语言和世界的知识。

通过这种``填空''任务，模型可以隐式地学到：

\begin{itemize}
    \item \textbf{世界知识}（trivia）：``Stanford University is located in \underline{Stanford/Palo Alto}, California''
    \item \textbf{句法}（syntax）：``I put \underline{the/a} fork down on the table''
    \item \textbf{共指}（coreference）：``The woman walked across the street, checking for traffic over \underline{her} shoulder''
    \item \textbf{语义}（semantics）：``I went to the ocean to see the fish, turtles, seals, and \underline{whales}''
    \item \textbf{情感}（sentiment）：``The movie was \underline{bad/great}''
    \item \textbf{推理}（reasoning）：``Iroh went into the kitchen... Zuko left the \underline{kitchen}''
\end{itemize}


% Raster figure omitted from this study copy.



\subsection{预训练--微调范式}


% Raster figure omitted from this study copy.



\begin{importantbox}{预训练--微调范式（Pretrain-Finetune Paradigm）}
\begin{enumerate}
    \item \textbf{预训练}：在海量无标注文本上训练，学习语言的通用表示。目标是最小化预训练损失 $\hat{\theta} = \arg\min_\theta \mathcal{L}_{\text{pretrain}}(\theta)$
    \item \textbf{微调}：以 $\hat{\theta}$ 为起点，在下游任务的标注数据上继续训练。梯度下降从 $\hat{\theta}$ 出发，走向对当前任务最优的参数 $\theta^*$
\end{enumerate}
\end{importantbox}

为什么预训练的参数初始化如此重要？Chris Manning 强调：``If you could actually solve this min, the starting point shouldn't matter. But it really, really, really does.'' 直觉上：

\begin{itemize}
    \item 梯度下降在微调过程中会停留在 $\hat{\theta}$ 附近的局部最优
    \item $\hat{\theta}$ 附近的参数空间已经``学会了语言''，因此该区域的局部最优往往\textbf{泛化能力更强}
    \item 无标注数据的量级（万亿词级别）远超标注数据（百万词级别），预训练利用了这个巨大的信息优势
\end{itemize}

\begin{warningbox}{预训练的数据优势不仅在于量}
即使你拥有海量的情感分析标注数据，从头训练的模型也可能不如预训练+微调的模型。原因在于：
\begin{itemize}
    \item 单一任务的标注数据多样性有限，模型容易过拟合到特定分布
    \item 预训练数据涵盖了互联网上极其多样的文本类型，使模型具备更强的\textbf{分布外泛化}（out-of-distribution generalization）能力
    \item 语言建模本身是一个极其困难的任务，要求模型学习语言的方方面面
\end{itemize}
\end{warningbox}

\subsection{本章小结}

从 Word2Vec 的``仅预训练嵌入''到``预训练整个网络''，NLP 经历了一次范式转变。预训练--微调范式的成功基于两个关键因素：（1）海量无标注文本提供了丰富的自监督信号；（2）预训练学到的参数提供了一个泛化能力极强的初始化点。这一范式至今仍是 NLP 的主流方法论。

%% ========== Part 3: Encoder 预训练 ==========

\section{Encoder 预训练：BERT 与掩码语言建模}

\subsection{三种 Transformer 架构回顾}

预训练方法与 Transformer 的架构类型密切相关。Chris Manning 将模型分为三类：


% Raster figure omitted from this study copy.



\begin{knowledgebox}{三种架构的特点与适用场景}
\begin{itemize}
    \item \textbf{Encoder}：每个 token 能看到整个序列（双向上下文）。适合\textbf{理解}类任务（分类、NER、问答抽取）。代表模型：BERT
    \item \textbf{Encoder-Decoder}：Encoder 部分看到完整输入，Decoder 部分自回归生成输出。适合\textbf{序列到序列}任务（翻译、摘要）。代表模型：T5
    \item \textbf{Decoder}：每个 token 只能看到它之前的 token（单向上下文）。适合\textbf{生成}类任务。代表模型：GPT 系列
\end{itemize}
\end{knowledgebox}

\subsection{为什么 Encoder 不能用语言建模预训练？}

Encoder 拥有双向上下文，这意味着在预测下一个词时，模型已经\textbf{看到了}那个词——预测任务变得毫无意义（trivial）。


% Raster figure omitted from this study copy.



解决方案：\textbf{掩码语言建模}（Masked Language Modeling, MLM）——先遮住一些词，再让模型预测被遮住的词。

\subsection{BERT：掩码语言建模}

BERT（Bidirectional Encoder Representations from Transformers，2018）是第一个成功将掩码语言建模大规模应用于预训练的模型。


% Raster figure omitted from this study copy.



BERT 的掩码策略选取 15\% 的 token 进行处理，对于被选中的 token：

\begin{itemize}
    \item \textbf{80\%} 替换为 \texttt{[MASK]} 标记
    \item \textbf{10\%} 替换为词表中的随机词
    \item \textbf{10\%} 保持不变
\end{itemize}

\begin{warningbox}{为什么不能全部用 [MASK]？}
如果只在 \texttt{[MASK]} 标记处做预测，模型可能会学到：``只有遇到 [MASK] 时才需要认真构建表示，其他位置可以偷懒''。但在实际使用时（如情感分析），输入中\textbf{不会出现} [MASK] 标记。通过混合使用随机替换和保持不变，模型被迫在\textbf{任何位置}都构建高质量的上下文表示。
\end{warningbox}


% Raster figure omitted from this study copy.

\footnotetext{Slides 第35页。来自 Devlin et al., 2018。}

\subsection{下一句预测任务（NSP）}

BERT 原始论文还引入了一个辅助任务——\textbf{下一句预测}（Next Sentence Prediction, NSP）：给定两段文本 A 和 B，预测 B 是否是 A 的真实续文。

这个任务的设计初衷是教会模型理解\textbf{长距离的文本连贯性}。然而，后续研究（RoBERTa，2019）发现这个任务\textbf{并不必要}，原因有二：

\begin{enumerate}
    \item NSP 将有效上下文长度减半（原本 512 token 的窗口被分成两段各 256 token）
    \item 模型在 NSP 任务上表现很差，可能是因为这个任务对当时的模型来说太难了
\end{enumerate}

\subsection{BERT 在 GLUE 基准上的表现}

BERT 的发布在 NLP 领域引起了巨大震动。在 GLUE 基准（包含多种 NLP 任务）上，BERT 全面超越了之前\textbf{为每个任务精心设计}的专用模型。


% Raster figure omitted from this study copy.

\footnotetext{Slides 第36页。来自 Devlin et al., 2018。}

Chris Manning 评价道：``All of that was blown out of the water by just build a big Transformer and teach it to predict the missing words.'' BERT 的成功宣告了``为每个任务设计专用架构''时代的终结。

\begin{importantbox}{BERT 的核心规格}
\begin{itemize}
    \item \textbf{BERT$_{\text{BASE}}$}：110M 参数，12 层，768 维隐层，12 头
    \item \textbf{BERT$_{\text{LARGE}}$}：340M 参数，24 层，1024 维隐层，16 头
    \item \textbf{训练数据}：BookCorpus + English Wikipedia（约 25 亿词）
    \item \textbf{预训练计算}：在当时被认为``只有 Google 才负担得起''，但今天看来并不算大
    \item \textbf{微调}：在单个 GPU 上即可完成
\end{itemize}
\end{importantbox}

\subsection{BERT 的微调方式}

将预训练好的 BERT 用于下游任务非常简单：


% Raster figure omitted from this study copy.



\begin{enumerate}
    \item 移除预训练时的 MLM 预测头
    \item 对于分类任务：取 \texttt{[CLS]} 标记或最后一个 token 的输出向量，接一个线性分类器
    \item 对于 token 级别任务（如 NER）：对每个 token 的输出向量接线性分类器
    \item 整个网络（BERT + 新分类器）一起微调
\end{enumerate}

\subsection{BERT 的局限性}

\begin{warningbox}{BERT 不适合文本生成}
BERT 的双向上下文使其非常擅长``理解''任务（分类、抽取式问答、NER），但它\textbf{没有自回归生成的能力}——无法像语言模型那样逐词生成文本。如果你需要生成摘要、翻译或对话回复，应该使用 Encoder-Decoder 或 Decoder-only 模型。
\end{warningbox}

\subsection{BERT 的改进：SpanBERT 与 RoBERTa}

\textbf{SpanBERT}：用\textbf{连续片段掩码}（span masking）替代随机单词掩码。随机掩码单个子词过于简单（前后子词的线索太强），而掩码一整个连续片段迫使模型进行更深层的理解。

\textbf{RoBERTa}（Robustly Optimized BERT Approach）：
\begin{itemize}
    \item 移除了 NSP 任务
    \item 在更多数据上训练更长时间
    \item 动态掩码（每次 epoch 重新生成掩码位置）
    \item 是 BERT 的\textbf{直接替代品}，性能更优
\end{itemize}

\begin{knowledgebox}{选择建议}
如果你的项目需要使用 encoder 模型（如分类、NER、抽取式问答），优先使用 \textbf{RoBERTa} 而非原始 BERT，它在几乎所有任务上都表现更好，且完全兼容 BERT 的使用方式。
\end{knowledgebox}

\subsection{参数高效微调}

当模型变得越来越大时，全参数微调变得昂贵且不实际。\textbf{参数高效微调}（Parameter-Efficient Fine-Tuning, PEFT）只调整少量参数，而冻结大部分预训练参数。


% Raster figure omitted from this study copy.



三种主要的参数高效微调方法：

\begin{enumerate}
    \item \textbf{Prefix Tuning / Prompt Tuning}：冻结网络，在序列前添加可训练的``伪词向量''
    \item \textbf{Adapter Layers}：在 Transformer 层之间插入小型可训练的适配层
    \item \textbf{LoRA}（Low-Rank Adaptation）：冻结原始权重矩阵 $W$，学习一个低秩增量 $\Delta W = BA$，令 $W' = W + BA$
\end{enumerate}

\begin{importantbox}{LoRA 的核心思想}
对于权重矩阵 $W \in \mathbb{R}^{d \times d}$，LoRA 学习两个小矩阵 $B \in \mathbb{R}^{d \times r}$ 和 $A \in \mathbb{R}^{r \times d}$（$r \ll d$），令实际权重为：
$$W' = W + BA$$
只训练 $B$ 和 $A$（参数量从 $d^2$ 降至 $2dr$），冻结 $W$。这背后的直觉是：从预训练到下游任务的参数变化是\textbf{低秩的}，不需要修改全部参数。
\end{importantbox}

\subsection{本章小结}

Encoder 预训练以 BERT 为代表，通过掩码语言建模让模型学会双向上下文表示。BERT 在各类 NLP 基准上的巨大成功标志着``预训练--微调''范式的确立。后续的 SpanBERT、RoBERTa 在训练策略上进行了优化。参数高效微调方法（Prefix Tuning、LoRA）则解决了大模型微调的效率问题。但 Encoder 模型不适合生成任务，这引出了 Encoder-Decoder 和 Decoder-only 架构。

%% ========== Part 4: Encoder-Decoder 预训练 ==========

\section{Encoder-Decoder 预训练：T5 与片段腐蚀}

\subsection{Encoder-Decoder 的预训练目标}

Encoder-Decoder 架构可以简单地用\textbf{语言建模}来预训练：将文本前半部分交给 Encoder（提供双向上下文），用 Decoder 自回归生成后半部分。


% Raster figure omitted from this study copy.

\footnotetext{Slides 第41页。来自 Raffel et al., 2018。}

但在实践中，一种更有效的方法是\textbf{片段腐蚀}（Span Corruption）。

\subsection{T5 与片段腐蚀}

T5（Text-to-Text Transfer Transformer）将所有 NLP 任务统一建模为``文本到文本''的格式，并使用片段腐蚀作为预训练目标。


% Raster figure omitted from this study copy.



片段腐蚀的工作方式：
\begin{enumerate}
    \item 在输入中随机遮住若干\textbf{连续片段}，用特殊标记（如 \texttt{[X]}、\texttt{[Y]}）替代
    \item Encoder 处理带遮罩的输入（保持双向上下文）
    \item Decoder 自回归\textbf{生成}被遮住的内容：先输出 \texttt{[X]}，再输出对应片段，再输出 \texttt{[Y]}，再输出对应片段...
\end{enumerate}

\begin{importantbox}{片段腐蚀 = BERT 的掩码思想 + 生成式输出}
片段腐蚀结合了两个世界的优点：
\begin{itemize}
    \item 像 BERT 一样利用\textbf{双向上下文}来理解输入
    \item 像语言模型一样以\textbf{自回归方式生成}输出
\end{itemize}
这使得预训练好的 T5 可以自然地应用于翻译、摘要等序列到序列任务。
\end{importantbox}

\subsection{T5 的``隐式知识检索''能力}

一个令人惊讶的发现是：T5 在预训练中看过大量类似``Franklin D. Roosevelt was born in [MASK]''的文本，之后经过在问答数据上微调，它能够\textbf{直接从参数中}回忆出答案，而非从外部知识库检索。

这种``闭卷''问答能力表明，预训练模型在其参数中\textbf{隐式存储了大量事实知识}。

\begin{warningbox}{模型回答总是看起来很流畅，但经常是错的}
Chris Manning 特别指出：这些模型的回答``always look very fluent, always look very reasonable, but they're frequently wrong.'' 这个问题（后来被称为``幻觉''，hallucination）至今仍是大语言模型的核心挑战之一。
\end{warningbox}

\subsection{本章小结}

Encoder-Decoder 预训练以 T5 为代表，通过片段腐蚀实现了双向理解与生成能力的统一。T5 将所有任务转化为``文本到文本''格式，提供了一种优雅的统一框架。模型在预训练过程中隐式存储了大量事实知识，但也伴随着幻觉问题。

%% ========== Part 5: Decoder 预训练（GPT 系列） ==========

\section{Decoder 预训练：GPT 系列与大语言模型}

\subsection{Decoder 的预训练与微调}

Decoder-only 模型的预训练目标就是标准的\textbf{语言建模}：

$$\mathcal{L}(\theta) = -\sum_{t=1}^{T} \log p_\theta(w_t | w_1, \ldots, w_{t-1})$$


% Raster figure omitted from this study copy.



微调方式与 BERT 类似：取序列最后一个 token 的隐状态 $h_T$，接一个线性层 $y \sim Ah_T + b$ 进行分类。

\subsection{GPT-1：开创性的预训练 Decoder}

\textbf{GPT}（Generative Pre-Training，2018，OpenAI）是第一个大规模预训练的 Decoder-only 模型：

\begin{itemize}
    \item 117M 参数，12 层，768 维隐层
    \item 约 40,000 词的 BPE 词表
    \item 在 BookCorpus 上训练
\end{itemize}


% Raster figure omitted from this study copy.



GPT 在多项 NLP 基准上取得了强劲表现。BERT 紧随其后发布，凭借双向上下文在分类任务上略胜一筹，但 GPT 的生成能力是 BERT 无法匹敌的。

\subsection{GPT-2：惊人的生成能力}

GPT-2（2019）将模型规模提升至 \textbf{1.5B 参数}，训练数据增至约 90 亿词。


% Raster figure omitted from this study copy.

\footnotetext{Slides 第51页。来自 Radford et al., 2019。}

GPT-2 的生成质量在当时令人震惊，它能够产生\textbf{多段落、主题连贯}的文本，且 1.5B 的规模仍然可以在小型 GPU 上进行微调。

\subsection{GPT-3：涌现的上下文学习能力}

GPT-3（2020）实现了参数规模的巨大飞跃：

\begin{table}[H]
\centering
\begin{tabular}{lccc}
\toprule
\textbf{模型} & \textbf{参数量} & \textbf{训练数据} & \textbf{年份} \\
\midrule
GPT-1 & 117M & BookCorpus & 2018 \\
GPT-2 & 1.5B & WebText（9B 词） & 2019 \\
GPT-3 & 175B & 300B 词 & 2020 \\
\bottomrule
\end{tabular}
\caption{GPT 系列模型的规模演进}
\end{table}

GPT-3 最令人惊讶的不是它的微调性能，而是它展现出了一种全新的能力——\textbf{上下文学习}（In-Context Learning, ICL）。


% Raster figure omitted from this study copy.

\footnotetext{Slides 第53页。来自 Brown et al., 2020。}

\begin{importantbox}{上下文学习（In-Context Learning）}
无需任何梯度更新或微调，仅通过在 prompt 中给出几个\textbf{示例}（demonstrations），模型就能``学会''执行新任务。这完全是通过\textbf{模式匹配}实现的：模型识别出 prompt 中的输入-输出模式，并将其应用于新的输入。

GPT-3 的上下文学习能力涵盖：
\begin{itemize}
    \item 翻译（给出几对翻译示例）
    \item 算术（给出几个加法示例）
    \item 纠错（给出几个错别字修正示例）
    \item 以及更多...
\end{itemize}
\end{importantbox}


% Raster figure omitted from this study copy.



Chris Manning 强调这种能力是``qualitatively new behavior''——从小模型中完全无法预见。GPT-3 只是被训练来预测下一个词，但它\textbf{涌现}出了执行复杂任务的能力。这既令人兴奋，也令人不安，因为我们无法准确预测什么能力会涌现、什么时候会涌现。

\subsection{Scaling Laws 与计算效率}


% Raster figure omitted from this study copy.

\footnotetext{Slides 第56页。来自 Kaplan et al., 2020。}

\begin{knowledgebox}{Chinchilla Scaling Laws}
DeepMind 的 Chinchilla 研究（Hoffmann et al., 2022）发现 GPT-3 的参数量与训练数据量之间\textbf{严重不平衡}——GPT-3 ``comically oversized''。Chinchilla 只有 70B 参数（不到 GPT-3 的一半），但在 1.4T token 上训练（GPT-3 只用了 300B），结果性能更优。

核心结论：给定固定计算预算，参数量 $N$ 和训练 token 数 $D$ 应该\textbf{等比例}缩放：$N \propto D$。大约每个参数对应 20 个训练 token。
\end{knowledgebox}

\subsection{Chain-of-Thought 提示}


% Raster figure omitted from this study copy.

\footnotetext{Slides 第59页。来自 Wei et al., 2022。}

\begin{importantbox}{Chain-of-Thought Prompting 的原理}
标准 prompting：问题 $\rightarrow$ 答案

Chain-of-Thought：问题 $\rightarrow$ \textbf{推理步骤} $\rightarrow$ 答案

关键洞察：模型在生成每个 token 时能\textbf{条件于之前生成的所有 token}。通过先生成中间推理步骤，模型：
\begin{itemize}
    \item 获得了额外的``思考空间''（类似草稿纸）
    \item 将复杂问题分解为更简单的子问题
    \item 在每个子步骤中能利用之前步骤的结论作为条件
\end{itemize}
这显著提升了数学推理、逻辑推理等需要多步思考的任务的准确率。
\end{importantbox}

\subsection{为什么 Decoder-only 模型占据主导地位？}

尽管 Encoder 和 Encoder-Decoder 各有优势，当今最大、最强的预训练模型几乎都是 Decoder-only 架构。Chris Manning 指出原因可能包括：

\begin{enumerate}
    \item \textbf{简洁性}：Decoder-only 架构更简单，所有参数共享在一个网络中
    \item \textbf{参数效率}：Encoder-Decoder 需要将参数分配给两个子网络，而 Decoder-only 的所有参数都用于同一个任务
    \item \textbf{统一性}：所有任务（理解+生成）都可以建模为``给定 prompt，生成 completion''
    \item \textbf{涌现能力}：上下文学习等涌现能力在 Decoder-only 模型中表现最为突出
\end{enumerate}

\subsection{本章小结}

Decoder-only 预训练以 GPT 系列为代表，通过简单的语言建模目标学习强大的语言能力。从 GPT-1（117M）到 GPT-3（175B），模型规模的增长带来了质变式的涌现能力，尤其是上下文学习。Chinchilla 揭示了``更大不一定更好''的计算效率规律。Chain-of-Thought prompting 则展示了引导模型进行多步推理的有效策略。Decoder-only 架构已成为当今大语言模型的主流选择。

%% ========== Part 6: 预训练数据 ==========

\section{预训练数据与伦理考量}

\subsection{预训练数据的来源}


% Raster figure omitted from this study copy.



主要的预训练数据来源包括：

\begin{itemize}
    \item \textbf{CommonCrawl}：互联网爬取数据（量最大，质量参差不齐）
    \item \textbf{Wikipedia}：高质量百科全书文本
    \item \textbf{BookCorpus / Books}：书籍语料
    \item \textbf{WebText / OpenWebText}：从 Reddit 高赞链接中收集的网页
    \item \textbf{学术论文}（ArXiv、PubMed）
    \item \textbf{代码}（GitHub、StackExchange）
\end{itemize}

\subsection{预训练学到了什么——以及它的阴暗面}


% Raster figure omitted from this study copy.



\begin{warningbox}{预训练的伦理风险}
Chris Manning 特别强调：模型在学习有用知识的同时，也会``learn and exacerbate racism and sexism, all manner of biases''。这些偏见来自训练数据中的不平衡和有害内容，而模型会忠实地复制甚至放大这些偏见。

这引出了关于预训练的重要伦理问题：
\begin{itemize}
    \item 训练数据的筛选和清洗至关重要
    \item 模型的输出不应被盲目信任
    \item 部署前需要仔细评估模型在不同群体上的表现差异
    \item 大模型的能力边界和失败模式仍不完全清楚
\end{itemize}
\end{warningbox}

\subsection{本章小结}

预训练数据的规模和多样性是模型能力的基础，但数据质量和伦理问题同样不可忽视。随着模型越来越大、能力越来越强，理解它们的工作原理和失败模式变得愈发重要——这不仅是为了用好它们，更是为了负责任地部署它们。

%% ========== 总结与延伸 ==========

\section{总结与延伸}

\subsection{核心知识图谱}

本节课建立了从词表示到大语言模型的完整认知链：

\begin{center}
\begin{tikzpicture}[node distance=0.9cm, every node/.style={draw, rounded corners, minimum width=4cm, minimum height=0.7cm, font=\small}]
\node (sub) {子词分词 (BPE)};
\node (ctx) [below=of sub] {上下文表示 vs 静态嵌入};
\node (para) [below=of ctx] {预训练--微调范式};
\node (enc) [below=of para] {Encoder: BERT (MLM)};
\node (encdec) [below=of enc] {Enc-Dec: T5 (Span Corruption)};
\node (dec) [below=of encdec] {Decoder: GPT (LM)};
\node (scale) [below=of dec] {规模化与涌现能力};
\draw[->, thick] (sub) -- (ctx);
\draw[->, thick] (ctx) -- (para);
\draw[->, thick] (para) -- (enc);
\draw[->, thick] (enc) -- (encdec);
\draw[->, thick] (encdec) -- (dec);
\draw[->, thick] (dec) -- (scale);
\node[draw=none, right=0.8cm of sub, text width=5cm, font=\scriptsize] {解决 OOV 问题，词表大小 vs 序列长度的权衡};
\node[draw=none, right=0.8cm of ctx, text width=5cm, font=\scriptsize] {Word2Vec 的一词一向量 vs Transformer 的上下文向量};
\node[draw=none, right=0.8cm of para, text width=5cm, font=\scriptsize] {海量无标注数据 $\rightarrow$ 少量标注数据微调};
\node[draw=none, right=0.8cm of enc, text width=5cm, font=\scriptsize] {双向上下文，擅长理解，不能生成};
\node[draw=none, right=0.8cm of encdec, text width=5cm, font=\scriptsize] {理解+生成，序列到序列任务};
\node[draw=none, right=0.8cm of dec, text width=5cm, font=\scriptsize] {单向上下文，擅长生成，规模化的主流选择};
\node[draw=none, right=0.8cm of scale, text width=5cm, font=\scriptsize] {ICL, CoT, Scaling Laws, Chinchilla};
\end{tikzpicture}
\end{center}

\subsection{关键 Takeaways}

\begin{importantbox}{本课七条核心原则}
\begin{enumerate}
    \item \textbf{子词分词解决了开放词表问题}：BPE 等算法使模型能够处理任意输入，消除了 UNK 标记
    \item \textbf{上下文表示优于静态嵌入}：同一个词在不同上下文中应有不同的向量表示
    \item \textbf{预训练--微调是当今 NLP 的基础范式}：先在海量无标注文本上学习通用表示，再用少量标注数据适配具体任务
    \item \textbf{三种架构各有所长}：Encoder（理解）、Encoder-Decoder（理解+生成）、Decoder（生成），但 Decoder-only 正成为主流
    \item \textbf{规模带来涌现能力}：上下文学习、Chain-of-Thought 推理等能力在大规模模型中``意外''出现
    \item \textbf{计算效率很重要}：Chinchilla 表明参数量与数据量应等比例缩放，盲目增大模型不一定最优
    \item \textbf{能力与风险并存}：预训练模型学到的知识包括有害偏见，部署时必须谨慎评估
\end{enumerate}
\end{importantbox}

\subsection{拓展阅读}

\begin{itemize}
    \item Devlin et al., \textbf{BERT: Pre-training of Deep Bidirectional Transformers for Language Understanding} (2018): \url{https://arxiv.org/abs/1810.04805}
    \item Liu et al., \textbf{RoBERTa: A Robustly Optimized BERT Pretraining Approach} (2019): \url{https://arxiv.org/abs/1907.11692}
    \item Joshi et al., \textbf{SpanBERT: Improving Pre-training by Representing and Predicting Spans} (2020): \url{https://arxiv.org/abs/1907.10529}
    \item Raffel et al., \textbf{Exploring the Limits of Transfer Learning with a Unified Text-to-Text Transformer} (T5, 2020): \url{https://arxiv.org/abs/1910.10683}
    \item Radford et al., \textbf{Improving Language Understanding by Generative Pre-Training} (GPT, 2018)
    \item Brown et al., \textbf{Language Models are Few-Shot Learners} (GPT-3, 2020): \url{https://arxiv.org/abs/2005.14165}
    \item Hoffmann et al., \textbf{Training Compute-Optimal Large Language Models} (Chinchilla, 2022): \url{https://arxiv.org/abs/2203.15556}
    \item Wei et al., \textbf{Chain-of-Thought Prompting Elicits Reasoning in Large Language Models} (2022): \url{https://arxiv.org/abs/2201.11903}
    \item Hu et al., \textbf{LoRA: Low-Rank Adaptation of Large Language Models} (2021): \url{https://arxiv.org/abs/2106.09685}
    \item Gururangan et al., \textbf{Don't Stop Pretraining} (2020): \url{https://arxiv.org/abs/2004.10964}
\end{itemize}

\end{document}
